\section{Prospective Real-Repository Construction Study}
\label{sec:real}

The final V4 work was a bounded construction-feasibility study, not an effect
estimation experiment.  Before target outcomes, the procedure fixed
artifact-bound target controls, top-one source locking, source correctness and
focal-safety evidence, a named \(p^\star\) applicability predicate, a
directional source--target shift, a no-second-major-mismatch review, memory
fidelity, resource controls, and a stopping rule.  Earlier protocol development
serves only as motivation for these final evidence gates; details are retained
in provenance artifacts rather than presented as a chronological repair diary.

The original five development targets produced
\VFourOriginalAcceptance{} complete acceptances.  Under the already frozen
trigger, a deterministic identity-only rule locked a five-target extension
before any selected target content was opened.  The extension produced
\VFourExtensionAcceptance{} complete acceptances and exhausted the authorized
development universe.  No second extension, relaxed gate, source fallback, or
confirmatory screening was permitted.

\begin{table*}[t]
\centering
\scriptsize
\caption{Gate-specific attrition in the bounded V4 development study. Counts characterize only the fixed development targets, not a population.}
\label{tab:v4-attrition}
\begin{tabularx}{\textwidth}{@{}llXXX@{}}
\toprule
Cohort & Target & Gate category & Terminal reason & Interpretation \\
\midrule
Original & aio-libs/aiohttp-session & SOURCE\_SIDE\_REJECTION & SOURCE\_SAFETY\_REJECT & Candidate ceased at this evidence gate \\
Original & buildbot/buildbot & SOURCE\_SIDE\_REJECTION & SOURCE\_SAFETY\_REJECT & Candidate ceased at this evidence gate \\
Original & wagtail/wagtail & TECHNICAL\_UNEVALUABILITY & TARGET\_TECHNICAL\_INVALID & No semantic inference \\
Original & django/django & SEMANTIC\_REJECTION & TASK\_STATEMENT\_CUE\_REJECT & Candidate ceased at this evidence gate \\
Original & psf/requests & SEMANTIC\_REJECTION & TASK\_STATEMENT\_CUE\_REJECT & Candidate ceased at this evidence gate \\
Locked extension & vyperlang/vyper & TARGET\_INVALIDITY & B\_U\_R\_TASK\_SECURITY\_MATRIX\_REJECT & Candidate ceased at this evidence gate \\
Locked extension & openstack/aodh & SOURCE\_SIDE\_REJECTION & NO\_SOURCE\_PASSES\_HARD\_GATES & Candidate ceased at this evidence gate \\
Locked extension & urllib3/urllib3 & SOURCE\_SIDE\_REJECTION & NO\_SOURCE\_PASSES\_HARD\_GATES & Candidate ceased at this evidence gate \\
Locked extension & tensorflow/tensorflow & TECHNICAL\_UNEVALUABILITY & TARGET\_TECHNICAL\_INVALID & Runtime/disk-quota failure during image materialization; no semantic inference \\
Locked extension & django/django & SOURCE\_SIDE\_REJECTION & NO\_SOURCE\_PASSES\_HARD\_GATES & Candidate ceased at this evidence gate \\
\bottomrule
\end{tabularx}
\vspace{2pt}
\parbox{.96\textwidth}{\footnotesize The original set yielded 0/5 complete acceptances; the prospectively locked extension yielded 0/5. TensorFlow remains \texttt{TECHNICAL\_UNDETERMINED}/\texttt{TARGET\_TECHNICAL\_INVALID}: a runtime quota failure is not semantic evidence. The resulting 0/10 is not a prevalence estimate.}
\end{table*}


Table~\ref{tab:v4-attrition} answers where plausible candidates ceased to
support the counterfactual: at task-statement semantics, source-side evidence,
the target task/security matrix, or technical evaluability.  TensorFlow's
runtime-quota failure remains
\texttt{TECHNICAL\_UNDETERMINED}/\texttt{TARGET\_TECHNICAL\_INVALID}; it is not
semantic rejection evidence.  Overall acceptance was
\VFourTotalAcceptance{}, under a small development-only sample.  This is not a
population estimate and supplies no inference about the frequency of valid
real-world pairs.

The frozen stop decision followed: no evaluated model call, GPU use, or unseen
confirmatory-target screening occurred.  Thus the study establishes
gate-specific construction attrition under its own bounded procedure, while
leaving the causal effect unestimated.
